臨床検査技師のための冊子ですが、第 1 部では検査ではなく、治療や予防などの介入の効果を統合するメタ解析を扱います。メタ解析の基本の手順は介入の研究をまとめる方法として作られ、検査のメタ解析はその手順の一部を変えたものだからです。第 1 部では、この基本の手順を、系統的レビューの工程の順にたどります。第 1 章でメタ解析が何を答え、何を答えられないかを確かめたあと、第 2 章で問いと計画を決め、第 3 章で文献を探して選び、第 4 章で論文から数字を取り出してバイアスのリスクを評価します。第 5 章では R で初めてのメタ解析を実行し、第 6 章で試験間のばらつきを、第 7 章で結果の確かさを調べます。第 8 章では、ここまでの作業を論文の形で報告する方法を説明します。例題は架空の降圧薬 A の 10 試験で、実データとして BCG ワクチンの 13 試験も使います。
